\section{Open Questions}

\begin{enumerate}
\item \textbf{Forward-QFT extension.}
Does the semiclassical construction extend cleanly to the forward QFT
(not just the inverse) with the same distribution-equivalence property,
and how does the feed-forward direction reverse?
\emph{Basis.} Griffiths and Niu present the semiclassical procedure for
the inverse QFT used in phase estimation. The forward QFT arises in
Shor-like subroutines and other algorithms where measurement is not the
terminal step, so the same substitution is not automatic. Our replication
only exercised the inverse direction.
\emph{Next steps.} Derive the forward-QFT analogue by conjugating the
iterative measure-and-feed-forward loop and verify distribution
equivalence against a coherent forward QFT applied to a battery of
computational-basis and superposition inputs at $k=3\text{--}5$, using the
same numpy harness.

\item \textbf{Integration with modern hybrid QPE schemes.}
How does the semiclassical iQFT integrate with modern hybrid schemes such
as iterative QPE (Kitaev), rejection-sampling QPE, and robust phase
estimation, and does it retain its exact-equivalence property when
combined with amplitude-amplification-based ancilla preparation?
\emph{Basis.} Since 1996 many phase-estimation variants have been proposed
that already assume measurement-based ancilla processing. Whether the
exact per-shot distribution equivalence still holds when the semiclassical
block is composed with these newer front-ends has not been checked here.
\emph{Next steps.} Pick two representative modern variants (e.g., Kitaev
iterative QPE and rejection-sampling QPE), simulate each with (a)
coherent iQFT tail and (b) semiclassical tail, and compare output
histograms across a grid of true phases. Quantify any deviation from
strict equivalence.

\item \textbf{Application to specific QPE benchmarks.}
Does the semiclassical iQFT provide a genuine circuit-runtime advantage on
realistic QPE benchmarks (e.g., electronic-structure Hamiltonians on the
H$_2$, LiH, or a spin-chain observable) once mid-circuit measurement
latency is included?
\emph{Basis.} This replication verified only algorithmic equivalence, not
runtime. Real hardware imposes a non-trivial time cost per mid-circuit
measurement (readout + classical feed-forward latency) that can offset the
saved two-qubit-gate time.
\emph{Next steps.} Build a QPE circuit for a small chemistry Hamiltonian
in Qiskit or Cirq, express it twice (all-quantum iQFT and semiclassical),
and use the noisy-simulator or hardware-calibration tables for a
superconducting device to compute total wall-clock and total
error-per-run for each variant.

\item \textbf{Noise robustness under measurement noise.}
How does the semiclassical iQFT degrade under realistic measurement noise,
in particular correlated readout errors and heralded misclassification,
compared to the coherent iQFT under equivalent gate noise?
\emph{Basis.} Griffiths and Niu assumed ideal projective measurement.
Under readout error, each fed-forward classical bit can be wrong,
propagating a coherent phase error to all subsequent qubits. This
asymmetry vs.\ coherent-iQFT gate errors was not exercised in our numerical
study.
\emph{Next steps.} Extend the numpy harness to inject an asymmetric
readout-error channel $(p_{0|1}, p_{1|0})$ and, in the coherent branch, an
equivalent depolarizing channel on the two-qubit gates. Sweep noise
strengths and plot phase-estimate bias/variance for both methods; compare
crossover points.

\item \textbf{Hardware-runtime overhead of mid-circuit measurement.}
What is the true hardware-runtime overhead of the mid-circuit measurement
and classical feed-forward step on current cloud QPU stacks, and does it
change the qualitative advantage of the semiclassical iQFT?
\emph{Basis.} The 1996 paper predates mid-circuit measurement being a
first-class primitive. On today's IBM Heron, IonQ Forte, and Quantinuum
H2 systems mid-circuit measurement has a specific, non-negligible cost
(microseconds on superconducting, milliseconds on trapped-ion). Whether
the semiclassical approach is faster end-to-end depends on this constant.
\emph{Next steps.} Query published latency numbers for mid-circuit
measurement plus conditional gate application on three current backends,
tabulate per-shot end-to-end time for $k=3\text{--}8$ QPE using coherent
vs.\ semiclassical iQFT, and identify the $k$ threshold at which each
backend crosses over.
\end{enumerate}
